% ============================================================================
% COURS 03 — EXERCICES DE LOIS ENTRÉE-SORTIE
% Source : 3-LoiES20222.docx
% ============================================================================

\section{Exercices — lois entrée-sortie}
\label{sec:ex-les}

\subsection{Système d'orientation d'antenne}

\textit{Source : Jérôme Letard — niveau 1.}

Un vérin électrique à vis-écrou, de pas $p=2\ \mathrm{mm/tr}$, oriente
une antenne parabolique. Le paramétrage utilise les points $A$, $B$ et
$C$, la longueur $AB=l_1$, la distance fixe $AC=l_0$ et les angles
$\alpha_1$ et $\alpha_2$.

\TLThreeImage[.60\linewidth]{images/image40.png}

\begin{enumerate}
  \item Identifier le paramètre d'entrée et le paramètre de sortie.
  \item Écrire la fermeture géométrique du triangle $ABC$.
  \item Déterminer la loi entrée-sortie reliant la longueur du vérin à l'orientation de l'antenne.
  \item Pour une variation angulaire imposée, calculer la variation de longueur puis le nombre de tours de la vis.
  \item En déduire la durée d'alimentation du moteur pour sa vitesse nominale.
\end{enumerate}

\subsection{Plan horizontal réglable de l'Airbus A340}

\textit{Source : CCP MP 2005 — niveau 2.}

Le plan horizontal réglable est commandé par une vis de manœuvre. Dans
le modèle simplifié, $O_4H=z$, $O_7H=a$ et :
\[
\overrightarrow{O_4O_7}=-a\vec x_b+b\vec z_b.
\]
Le PHR pivote d'un angle $\beta$ autour de $O_7$.

\begin{center}
\TLThreeImage[.43\linewidth]{images/image58.jpeg}\hfill
\TLThreeImage[.50\linewidth]{images/image66.png}
\end{center}

\begin{enumerate}
  \item Proposer une technologie limitant le jeu entre la vis et l'écrou.
  \item Compléter le modèle cinématique simplifié.
  \item Dessiner les figures de changement de base.
  \item Établir la loi $z=f(\beta)$.
  \item Tracer $\Delta z=z(\beta)-z(0)$ pour $-12^\circ\leq\beta\leq12^\circ$, avec $a=2\ \mathrm m$ et $b=0{,}5\ \mathrm m$.
  \item Justifier une approximation linéaire $\beta=K_v\Delta z$ et identifier $K_v$.
  \item Déterminer la longueur utile de la vis pour $-12^\circ\leq\beta\leq4^\circ$.
\end{enumerate}

\subsection{Échelle pivotante automatique EPAS}

\textit{Source : Centrale-Supélec PSI 2007 — niveau 1.}

Pendant le dressage, un vérin articulé en $B$ et $C$ fait pivoter le
berceau autour de $A$. On donne :
\[
O_0A=a,\quad AC=c,\quad O_0B=b,\quad BC=r(t),
\]
avec $a=2\ \mathrm m$, $b=3\ \mathrm m$ et $c=2{,}6\ \mathrm m$.

\TLEPASGeometry

\begin{enumerate}
  \item Écrire une fermeture géométrique adaptée.
  \item Établir $r=f(\theta)$.
  \item Calculer la course minimale entre $\theta=10^\circ$ et $\theta=50^\circ$.
  \item Vérifier que le sens de variation de $r$ est cohérent avec le mouvement du vérin.
\end{enumerate}

\subsection{Stabilisateur cardiaque actif Cardiolock}

\textit{Source : Centrale-Supélec TSI 2012 — niveau 3.}

Le mécanisme emploie des pivots élastiques sans jeu et un actionneur
piézoélectrique. Les petits déplacements permettent les approximations
$\sin\alpha\simeq\alpha$ et $\cos\alpha\simeq1$.

\begin{center}
\TLThreeImage[.31\linewidth]{images/image83.png}\hfill
\TLThreeImage[.31\linewidth]{images/image84.png}\hfill
\TLThreeImage[.31\linewidth]{images/image85.png}
\end{center}

\begin{enumerate}
  \item Écrire la fermeture de la chaîne $0-1-2-3-4-0$.
  \item Projeter et linéariser les équations reliant $\lambda$, $\alpha$ et $\gamma$.
  \item Déterminer $\lambda_0$ dans la configuration de repos, puis $x_A=\lambda-\lambda_0$.
  \item Pour $x_{A\max}=-130\ \mathrm{\mu m}$, calculer $\alpha_{\max}$ et $\gamma_{\max}$.
  \item Exprimer le déplacement vertical du point $C$, puis celui de l'extrémité $E$.
  \item Conclure sur la capacité à compenser le déplacement résiduel du cœur.
\end{enumerate}

\subsection{Compacteur Caterpillar}

\textit{Source : Centrale-Supélec 2004 — niveau 2.}

Deux demi-bâtis sont articulés et orientés par un vérin hydraulique. Le
cahier des charges impose un angle relatif compris entre $-32^\circ$ et
$+32^\circ$.

\TLThreeImage[.70\linewidth]{images/image90.png}

\begin{enumerate}
  \item Identifier les classes d'équivalence cinématique.
  \item Construire le graphe des liaisons.
  \item Choisir un paramétrage géométrique complet.
  \item Écrire la fermeture géométrique et établir la loi entrée-sortie.
  \item En déduire la course minimale du vérin.
\end{enumerate}

\subsection{Barrière Sinusmatic}

\textit{Niveau 3.}

Le dispositif transforme la rotation continue du plateau 7 en une
rotation alternative d'un quart de tour de la barrière 2. Le croisillon
3 impose une orthogonalité géométrique et un angle constant
$\gamma=45^\circ$.

\begin{center}
\TLThreeImage[.44\linewidth]{images/image98.png}\hfill
\TLThreeImage[.44\linewidth]{images/image100.png}
\end{center}

\begin{enumerate}
  \item Identifier les paramètres d'entrée $\alpha$ et de sortie $\beta$.
  \item Dessiner le graphe des liaisons.
  \item Représenter les figures de changement de base.
  \item Traduire l'orthogonalité $\vec y_3\perp\vec x_2$ par un produit scalaire.
  \item Déterminer la loi entrée-sortie en position.
  \item Vérifier l'amplitude de rotation et l'irréversibilité du mécanisme.
\end{enumerate}

\subsection{Mécanisme oscillant}

\TLThreeImage[.66\linewidth]{images/image105.png}

\begin{enumerate}
  \item Identifier la chaîne fermée et les paramètres de position.
  \item Écrire la fermeture géométrique.
  \item Éliminer le paramètre intermédiaire et établir la loi de l'angle de sortie.
  \item Déterminer les positions extrêmes et l'amplitude d'oscillation.
\end{enumerate}

\subsection{Mécanisme à came}

\TLThreeImage[.58\linewidth]{images/image106.png}

\begin{enumerate}
  \item Identifier le mouvement d'entrée et le mouvement de sortie.
  \item Définir la levée du poussoir en fonction de l'angle de came.
  \item Déterminer les phases de montée, de maintien et de descente.
  \item Déduire la vitesse du poussoir par dérivation de la loi de levée.
\end{enumerate}

\subsection{Bras de robot}

\begin{center}
\TLThreeImage[.46\linewidth]{images/image107.png}\hfill
\TLThreeImage[.46\linewidth]{images/image108.png}
\end{center}

\begin{enumerate}
  \item Identifier les coordonnées articulaires et opérationnelles.
  \item Écrire le vecteur position de l'effecteur par Chasles.
  \item Établir le modèle géométrique direct.
  \item Rechercher les configurations permettant d'atteindre un point imposé.
  \item Décrire les singularités et les limites de l'espace de travail.
\end{enumerate}

\subsection{Pompe oscillante}

\begin{center}
\TLThreeImage[.31\linewidth]{images/image109.png}\hfill
\TLThreeImage[.31\linewidth]{images/image110.png}\hfill
\TLThreeImage[.31\linewidth]{images/image111.png}
\end{center}

\begin{enumerate}
  \item Repérer l'excentrique, la bielle et le solide oscillant.
  \item Écrire la fermeture géométrique.
  \item Déterminer la loi en position puis la loi en vitesse.
  \item Identifier la course utile et les positions de point mort.
\end{enumerate}

\subsection{Banc d'épreuve hydraulique}

\begin{center}
\TLThreeImage[.22\linewidth]{images/image112.png}\hfill
\TLThreeImage[.22\linewidth]{images/image113.png}\hfill
\TLThreeImage[.22\linewidth]{images/image114.png}\hfill
\TLThreeImage[.22\linewidth]{images/image115.png}
\end{center}

\begin{enumerate}
  \item Construire le schéma cinématique à partir du dessin d'ensemble.
  \item Identifier entrée, sortie et paramètres intermédiaires.
  \item Établir la loi géométrique du mécanisme.
  \item Déterminer la course de l'actionneur nécessaire à l'effort d'épreuve.
\end{enumerate}

\subsection{Agitateur de cellules pancréatiques}

\textit{Source : CCP PSI 2010 — niveau 3.}

La chaîne principale associe un excentrique 1, une bielle 2 et un bras 3.
On donne :
\[
O_1O_2=e,\quad O_2B=b,\quad O_3B=L,
\quad \overrightarrow{O_3O_1}=c\vec x-d\vec y,
\quad \overrightarrow{O_3A}=q\vec x_3,
\]
avec $q=400\ \mathrm{mm}$.

\TLThreeImage[.64\linewidth]{images/image117.png}

\begin{enumerate}
  \item Écrire la fermeture géométrique et les deux équations projetées liant $\theta_1$, $\theta_2$ et $\theta_3$.
  \item Dériver les équations et déterminer la loi entrée-sortie en vitesse.
  \item Exploiter :
  \[
  \frac{\dot\theta_3}{\dot\theta_1}
  =\frac{e}{L}
  \frac{L\sin(\theta_1-\theta_3)-c\sin\theta_1-d\cos\theta_1}
  {e\sin(\theta_1-\theta_3)-d\cos\theta_3-c\sin\theta_3}.
  \]
  \item Qualifier la transformation de mouvement et discuter l'homocinétisme.
  \item Déterminer la loi en position et le domaine de variation de $\theta_3$.
  \item Exprimer la projection verticale de $\overrightarrow{O_3A}$ et vérifier la course annoncée de $100\ \mathrm{mm}$.
\end{enumerate}
